% =======================
% Section 5: Demo
% =======================
\section{Demo}

% ---- Slide 15: Transition to the Notebook Demo ----
\begin{frame}{Transition to the Notebook Demo}
\begin{columns}[T]
\begin{column}{0.52\textwidth}
\textbf{Next:} build a simple agent in AWS\\[0.6em]
Goals of the demo:
\begin{itemize}
  \item Define tools (and expose via MCP)
  \item Configure the model
  \item Create the agent
  \item Interact with it
\end{itemize}
\vspace{0.6em}
Focus: making the abstractions \textbf{concrete}\\[0.4em]
\textit{``Now that we have the conceptual map, we move to a practical implementation.''}
\end{column}
\begin{column}{0.44\textwidth}
\begin{tikzpicture}[
  box/.style={draw, rounded corners=4pt, minimum width=3.0cm,
              minimum height=0.7cm, align=center, font=\small},
  every node/.style={font=\small}
]
  \node[box, fill=blue!12]  (concept) at (0,1.6) {Concepts};
  \node[box, fill=green!20] (impl)    at (0,0.4) {Implementation};
  \draw[-{Stealth}, very thick] (concept) -- (impl);
\end{tikzpicture}
\end{column}
\end{columns}
% Speaker notes:
% The notebook is intentionally simpler than the full theory stack.
% Not every production concern will be implemented live,
% but students will see the essential agent loop.
\end{frame}

% =======================
% Appendix
% =======================
\appendix

\begin{frame}{Appendix A — Memory Types}
\begin{itemize}
  \item \textbf{Working / context memory} — current prompt and conversation window; ephemeral
  \item \textbf{Short-term / session memory} — state maintained across turns within a session
  \item \textbf{Long-term / persistent memory} — durable storage across sessions (external store)
\end{itemize}
\vspace{0.4em}
\small \textit{Memory captures what the agent experienced. The knowledge base captures what exists in the domain. Both feed into the context window at inference time — they are complementary, not alternatives.}
\end{frame}

\begin{frame}{Appendix B — Failure Modes of Agents}
\begin{itemize}
  \item \textbf{Hallucination} — confident but incorrect outputs
  \item \textbf{Wrong tool selection} — choosing an inappropriate tool for a task
  \item \textbf{Looping} — repeating the same steps without progress
  \item \textbf{Stale retrieval} — outdated or irrelevant chunks returned by RAG
  \item \textbf{Prompt injection} — malicious instructions embedded in retrieved content or tool outputs hijack the agent
  \item \textbf{Permission leakage} — agent accesses data or tools beyond its authorization
  \item \textbf{Excessive cost / latency} — runaway token usage or slow chains; the growing context window is a key structural driver
\end{itemize}
\end{frame}

\begin{frame}{Appendix C — Alternative Agent Execution Patterns}
\begin{itemize}
  \item \textbf{ReAct} — Thought $\to$ Act $\to$ Observe loop; grounded in real-world feedback at each step
  \item \textbf{ReWOO} — plans all tool calls upfront without an observation loop; more efficient but less adaptive
  \item \textbf{Reflexion} — adds iterative self-reflection; agent identifies errors and revises strategy
  \item \textbf{LATS} (Language Agent Tree Search) — builds a decision tree over possible actions using Monte Carlo search + self-reflection; excels at complex coding and QA tasks
\end{itemize}
\vspace{0.4em}
\small All of these are single-agent execution patterns. They operate at a different level from orchestration architecture (single-agent vs.\ multi-agent system design).
\end{frame}
